% ======================================================================
% capitulos/linea2/capitulo10.tex
% ----------------------------------------------------------------------
% LÍNEA CE02 — Evidencias CE0221 (modelo de datos relacional documentado y
% validado) y CE0222 (base de datos relacional implementada y consultable)
% Fuente: backend/app/models.py y backend/alembic/versions (rama main)
% Scripts: sql/01_ddl.sql, sql/02_dml.sql, sql/03_consultas.sql
% ======================================================================
\chapter{Plataforma de Datos del Sistema}
\label{chap:ce02-datos}

\entregable{CE02 --- Ingeniería de Software · competencia CE022 Ingeniería de la Información}{Semestre 2 · Entregable 2: Plataforma de Datos del Sistema}{CE0221 Modelo relacional · CE0222 BD implementada · CE0223 Motor transaccional programado · CE0224 BD administrada, asegurada y operativa}{Modelo y BD del producto real; programación, seguridad y respaldo probados en laboratorio, pendientes de pasar a producción.}


\section*{Resumen ejecutivo}
\addcontentsline{toc}{section}{Resumen ejecutivo}

Los datos son el centro del sistema de WISP: de la base de datos salen las credenciales
PPPoE que se crean en los routers, las fechas que deciden cada corte automático, el
mapa de bornes que usa el técnico en campo y los ingresos que revisa la gerencia. Una
inconsistencia en la base se traduce en un cliente cortado por error o en un borne
asignado dos veces.

Este entregable documenta la plataforma de datos completa: el \textbf{modelo} (21
tablas, 16 tipos enumerados, diccionario y análisis de normalización), la
\textbf{implementación} versionada con 32 migraciones Alembic y probada en
PostgreSQL~15, las \textbf{consultas y la programación} en la base (procedimiento de
registro de pagos con reactivación atómica, triggers de auditoría y de reglas de planta,
funciones e índices medidos con planes de ejecución) y la \textbf{seguridad y
administración} (roles con mínimo privilegio, auditoría inalterable, respaldo y
restauración probados, y consultas de monitoreo).

Todas las evidencias se ejecutaron sobre el esquema real del sistema, en un contenedor
\texttt{postgres:15-alpine} limpio y con datos anonimizados. Los resultados más
relevantes: los índices propuestos reducen la consulta del historial de pagos de
21,0~ms a 0,09~ms sobre 600\,000 pagos, y una restauración completa de 650\,000 filas
tomó unos 4~segundos con conteos y sumas idénticos a la base original. La programación
en base de datos y los roles están probados en laboratorio y pendientes de pasar a
producción mediante una migración Alembic.

% ======================================================================
\section{Sección 1: Modelo de datos}
\label{sec:ce02-modelo}

\subsection{Modelo entidad-relación}
\label{sec:ce02-er}

La base de datos tiene \textbf{21 tablas} y \textbf{16 tipos enumerados}. La
\cref{fig:ce02-er} muestra las entidades y relaciones principales.

\tikzset{ent/.style={draw=upeublue,fill=upeublue!6,font=\tiny,align=left,
  text width=2.75cm,inner sep=2.5pt,rounded corners=1pt},
  rel/.style={draw=upeublue!70,thick},
  card/.style={font=\tiny\bfseries,text=red!60!black,inner sep=1pt}}

\begin{figure}[htbp]
\centering
\begin{tikzpicture}[x=1cm,y=1cm,scale=0.82,every node/.style={transform shape}]
  % Fila 0
  \node[ent] (usu) at (0,0) {\textbf{usuarios}\\PK id · email (UQ), rol};
  \node[ent] (tz)  at (3.5,0) {\textbf{tecnico\_zona}\\FK tecnico\_id, zona\_id\\UQ(tecnico\_id, zona\_id)};
  \node[ent] (zon) at (7,0) {\textbf{zonas}\\PK id · nombre (UQ)\\ip\_wireguard, lista\_corte};
  \node[ent] (muf) at (10.5,0) {\textbf{mufas}\\FK zona\_id · codigo (UQ)};
  \node[ent] (cor) at (14,0) {\textbf{corridas}\\FK mufa\_id · metros};
  % Fila 1
  \node[ent] (jor) at (0,-2.7) {\textbf{jornadas\_vendedor}\\FK vendedor\_id, zona\_id};
  \node[ent] (gas) at (3.5,-2.7) {\textbf{gastos}\\FK zona\_id · monto};
  \node[ent] (caj) at (10.5,-2.7) {\textbf{cajas\_nap}\\FK zona\_id, mufa\_id\\codigo (UQ)};
  \node[ent] (ret) at (14,-2.7) {\textbf{ordenes\_retiro}\\FK cliente\_id, onu\_retirada\_id};
  % Fila 2
  \node[ent] (ubi) at (0,-5.4) {\textbf{ubicaciones\_vendedor}\\FK jornada\_id · lat, lng};
  \node[ent] (pla) at (3.5,-5.4) {\textbf{planes}\\PK id · codigo (UQ)};
  \node[ent,fill=upeublue!15] (cli) at (7,-5.4) {\textbf{clientes}\\FK zona\_id, plan\_id\\dni (UQ), usuario\_pppoe (UQ)\\estado, fecha\_vencimiento};
  \node[ent] (bor) at (10.5,-5.4) {\textbf{bornes}\\FK caja\_nap\_id, cliente\_id\\UQ(caja\_nap\_id, numero)};
  % Fila 3
  \node[ent] (pag) at (3.5,-8.1) {\textbf{pagos}\\FK cliente\_id\\numero\_operacion (UQ)};
  \node[ent] (his) at (7,-8.1) {\textbf{historial}\\FK cliente\_id, usuario\_id};
  \node[ent] (ave) at (10.5,-8.1) {\textbf{averias}\\FK cliente\_id, caja\_nap\_id,\\borne\_id, onu\_*};
  \node[ent] (onu) at (14,-8.1) {\textbf{onus}\\FK tecnico\_id, cliente\_id\\numero\_serie (UQ)};
  % Fila 4
  \node[ent] (mm) at (0,-10.8) {\textbf{mensajes\_masivos}\\FK creado\_por\_id};
  \node[ent] (mme) at (3.5,-10.8) {\textbf{mensajes\_masivos\_envios}\\FK mensaje\_masivo\_id, cliente\_id};
  \node[ent] (evi) at (10.5,-10.8) {\textbf{evidencias}\\FK averia\_id, cliente\_id,\\orden\_retiro\_id};
  % Relaciones (1 ... N)
  \newcommand{\relac}[3]{\draw[rel] (#1) -- node[card,pos=0.15]{1} node[card,pos=0.85]{#3} (#2);}
  \relac{usu}{tz}{N} \relac{zon}{tz}{N} \relac{zon}{muf}{N} \relac{muf}{cor}{N}
  \relac{zon}{cli}{N} \relac{pla}{cli}{N} \relac{zon}{caj}{N} \relac{muf}{caj}{N}
  \relac{caj}{bor}{N} \relac{cli}{bor}{0..1} \relac{cli}{pag}{N} \relac{cli}{his}{N}
  \relac{cli}{onu}{N} \relac{cli}{ave}{N} \relac{ave}{evi}{N} \relac{cli}{ret}{N}
  \relac{usu}{jor}{N} \relac{jor}{ubi}{N} \relac{zon}{gas}{N} \relac{mm}{mme}{N}
  \relac{cli}{mme}{N}
\end{tikzpicture}
\caption{Modelo entidad-relación (tablas principales; cardinalidad 1:N salvo indicación).
Además, \texttt{clientes} referencia a \texttt{usuarios} tres veces (vendedor, aprobador e
instalador), \texttt{onus} y \texttt{averias} referencian al técnico, y la tabla
\texttt{consultas\_documento\_mes} es independiente.}
\label{fig:ce02-er}
\end{figure}

% ======================================================================
\subsection{Modelo lógico relacional}
\label{sec:ce02-logico}

{\small
\begin{longtable}{L{3.4cm}L{5.6cm}L{5.0cm}}
\caption{Modelo lógico: tablas, claves foráneas y restricciones de unicidad.}\label{tab:ce02-logico}\\
\toprule
\textbf{Tabla} & \textbf{Claves foráneas} & \textbf{Unicidad / índices} \\
\midrule
\endfirsthead
\toprule
\textbf{Tabla} & \textbf{Claves foráneas} & \textbf{Unicidad / índices} \\
\midrule
\endhead
\bottomrule
\endlastfoot
usuarios & --- & \texttt{email} \\
tecnico\_zona & tecnico\_id $\rightarrow$ usuarios, zona\_id $\rightarrow$ zonas & (tecnico\_id, zona\_id) \\
zonas & --- & \texttt{nombre} \\
mufas & zona\_id $\rightarrow$ zonas & \texttt{codigo} \\
corridas & mufa\_id $\rightarrow$ mufas & --- \\
cajas\_nap & zona\_id $\rightarrow$ zonas, mufa\_id $\rightarrow$ mufas & \texttt{codigo} \\
bornes & caja\_nap\_id $\rightarrow$ cajas\_nap, cliente\_id $\rightarrow$ clientes & (caja\_nap\_id, numero) \\
planes & --- & \texttt{codigo} \\
clientes & zona\_id, plan\_id, vendedor\_id, aprobado\_por\_id, tecnico\_instalador\_id & \texttt{dni}, \texttt{ruc}, \texttt{usuario\_pppoe}; índice en \texttt{estado}, \texttt{migrado} \\
onus & tecnico\_id $\rightarrow$ usuarios, cliente\_id $\rightarrow$ clientes & \texttt{numero\_serie}; índice en \texttt{estado} \\
ordenes\_retiro & cliente\_id, tecnico\_asignado\_id, onu\_retirada\_id & índice en \texttt{estado} \\
averias & cliente\_id, caja\_nap\_id, borne\_id, tecnico\_asignado\_id, onu\_reemplazada\_id, onu\_nueva\_id & índice en \texttt{estado} \\
evidencias & averia\_id, cliente\_id, orden\_retiro\_id, subido\_por\_id & --- \\
pagos & cliente\_id $\rightarrow$ clientes, registrado\_por $\rightarrow$ usuarios & \texttt{numero\_operacion}; índice en \texttt{grupo\_adelanto} \\
gastos & zona\_id, registrado\_por & --- \\
historial & cliente\_id, usuario\_id & --- \\
mensajes\_masivos & creado\_por\_id $\rightarrow$ usuarios & índice en \texttt{estado} \\
mensajes\_masivos\_\allowbreak envios & mensaje\_masivo\_id, cliente\_id & índice en \texttt{estado} \\
jornadas\_vendedor & vendedor\_id, zona\_id & índice en \texttt{estado} \\
ubicaciones\_vendedor & jornada\_id $\rightarrow$ jornadas\_vendedor & --- \\
consultas\_documento\_mes & --- & (anio, mes) \\
\end{longtable}
}

% ======================================================================
\subsection{Diccionario de datos}
\label{sec:ce02-diccionario}

Se documentan las tablas centrales del negocio. Convenciones: \textbf{PK} clave
primaria, \textbf{FK} clave foránea, \textbf{UQ} único, \textbf{NN} no nulo.

\newcommand{\dictab}[2]{\multicolumn{4}{l}{\cellcolor{upeublue!10}\textbf{#1} --- \textit{#2}}\\}

{\small
\begin{longtable}{L{3.4cm}L{2.5cm}L{1.4cm}L{6.4cm}}
\caption{Diccionario de datos de las tablas centrales.}\label{tab:ce02-diccionario}\\
\toprule
\textbf{Campo} & \textbf{Tipo} & \textbf{Restr.} & \textbf{Descripción} \\
\midrule
\endfirsthead
\toprule
\textbf{Campo} & \textbf{Tipo} & \textbf{Restr.} & \textbf{Descripción} \\
\midrule
\endhead
\midrule \multicolumn{4}{r}{\textit{Continúa en la página siguiente}} \\
\endfoot
\bottomrule
\endlastfoot
\dictab{usuarios}{personal de la empresa con acceso al sistema}
id & integer & PK & Identificador. \\
nombre & varchar(100) & NN & Nombre completo. \\
email & varchar(150) & UQ, NN & Usuario de inicio de sesión. \\
hashed\_password & varchar(255) & NN & Hash bcrypt de la contraseña. \\
rol & rolusuario & NN & admin, tecnico, vendedor\_campo, atencion\_cliente. \\
activo & boolean & & Permite deshabilitar sin borrar. \\
telefono, dni, direccion, fecha\_ingreso & varios & & Ficha de personal (informativa). \\
\dictab{zonas}{zona de servicio con su router MikroTik}
id & integer & PK & Identificador. \\
nombre & varchar(100) & UQ, NN & Nombre de la zona (p.\,ej. SUCASCO). \\
ip\_lan & varchar(20) & & Gateway de la LAN de clientes de la zona. \\
ip\_wireguard & varchar(20) & & IP del router en el túnel (10.10.0.0/24); vacía si aún no hay túnel. \\
wg\_public\_key & varchar(255) & & Llave pública WireGuard del router. \\
pool\_pppoe & varchar(100) & & Nombre del pool de IP del MikroTik para los perfiles. \\
lista\_corte & varchar(100) & & Nombre exacto de la address-list de corte de esa zona. \\
usuario\_mikrotik, password\_mikrotik & varchar & & Credenciales heredadas (ver \cref{sec:ce02-normalizacion}). \\
lat, lng, descripcion, tiendas\_pago & varios & & Ubicación e información pública. \\
\dictab{cajas\_nap y bornes}{planta externa de fibra}
cajas\_nap.codigo & varchar(50) & UQ, NN & Código de la caja. \\
cajas\_nap.zona\_id / mufa\_id & integer & FK & Zona y mufa de la que cuelga. \\
cajas\_nap.capacidad\_bornes & integer & NN & Bornes de la caja (8 por defecto). \\
bornes.caja\_nap\_id, numero & integer & FK, NN & Caja y número; UQ(caja, número). \\
bornes.estado & estadoborne & NN & libre, ocupado, danado. \\
bornes.cliente\_id & integer & FK & Cliente conectado (nulo si libre). \\
bornes.codigo\_precinto & varchar(50) & & Precinto de seguridad colocado. \\
bornes.potencia\_optica\_dbm & float & & Potencia medida en la instalación. \\
\dictab{planes}{planes comerciales}
codigo & varchar(20) & UQ, NN & Código y nombre del perfil PPPoE (PLAN01\ldots). \\
bajada\_mbps, subida\_mbps & integer & NN & Velocidades contratadas. \\
precio & float & NN & Precio mensual (S/). \\
incluye\_tv, tv\_incluidas, precio\_tv\_adicional & varios & & Servicio de TV. \\
tipo\_conexion & tipoconexion & NN & fibra\_optica, inalambrico, satelital. \\
costo\_proveedor & float & & Solo satelital: costo pagado al proveedor. \\
\dictab{clientes}{abonados del servicio}
id & integer & PK & Identificador. \\
nombre & varchar(150) & NN & Nombre o razón social. \\
tipo\_documento & tipodocumento & NN & dni o ruc. \\
dni / ruc & varchar(15) / (11) & UQ & Documento; del DNI salen las credenciales PPPoE. \\
telefono\_whatsapp & varchar(20) & & Destino de los avisos. \\
lat, lng, direccion & varios & & Ubicación de la vivienda (GPS). \\
usuario\_pppoe & varchar(100) & UQ, NN & Usuario PPPoE en el MikroTik. \\
password\_pppoe & varchar(100) & NN & Contraseña PPPoE (ver \cref{sec:ce02-normalizacion}). \\
zona\_id & integer & FK & Zona (nula solo para satelitales). \\
plan\_id & integer & FK, NN & Plan contratado. \\
fecha\_instalacion, fecha\_vencimiento & date & & Ciclo de facturación de 30 días. \\
pausa\_hasta, compromiso\_pago\_hasta & date & & Pausa del servicio y prórroga acordada. \\
estado & estadocliente & NN & Ciclo de vida (\cref{fig:ce02-est-cliente}). \\
migrado & boolean & NN & Cliente anterior al sistema (sin datos de planta). \\
vendedor\_id, aprobado\_por\_id, tecnico\_instalador\_id & integer & FK & Responsables de cada etapa. \\
pin\_app & varchar(255) & & Hash del PIN de la App del cliente. \\
equipo\_* & varios & & Datos de equipo heredados (reemplazados por \texttt{onus}). \\
\dictab{onus}{equipos terminales del cliente}
numero\_serie & varchar(100) & UQ, NN & Serie del equipo (escaneada). \\
mac, marca, wifi\_version & varios & & Datos técnicos (wifi5, wifi6). \\
estado & estadoonu & NN & Ciclo de vida (\cref{fig:ce02-est-onu}). \\
tecnico\_id / cliente\_id & integer & FK & Técnico que lo tiene / cliente donde está instalado. \\
\dictab{pagos}{cobros mensuales}
cliente\_id & integer & FK, NN & Cliente que paga. \\
monto & float & NN & Importe (S/). \\
mes\_pagado & varchar(7) & NN & Mes cubierto (AAAA-MM). \\
fecha\_pago & date & NN & Fecha del cobro. \\
metodo\_pago & varchar(50) & & efectivo, yape, transferencia\ldots \\
numero\_operacion & varchar(100) & UQ & Operación bancaria/Yape (RN-06). \\
grupo\_adelanto, descuento & varchar(36), float & & Adelantos de varios meses con un solo recibo. \\
registrado\_por & integer & FK & Usuario que registró el pago. \\
\dictab{historial}{bitácora de eventos del cliente (RNF-08)}
cliente\_id, usuario\_id & integer & FK & Cliente afectado y autor. \\
tipo & tipohistorial & NN & activacion, corte, pago, cambio\_plan, cambio\_onu, compromiso\_pago, nota. \\
descripcion, creado\_en & text, timestamptz & NN & Detalle y momento del evento. \\
\end{longtable}
}

% ======================================================================
\subsection{Normalización y calidad del modelo}
\label{sec:ce02-normalizacion}

\begin{itemize}
  \item \textbf{1FN:} todas las tablas tienen clave primaria y atributos atómicos,
  salvo las columnas JSON \texttt{mufas.trazado}, \texttt{corridas.trazado} y
  \texttt{mensajes\_masivos.zona\_ids/estados\_incluidos}. Es una excepción
  consciente: son listas que solo se dibujan o se leen enteras y nunca se filtran
  por un elemento.
  \item \textbf{2FN:} todas las claves primarias son simples (\texttt{id}), así que no
  hay dependencias parciales; las combinaciones naturales se protegen con
  restricciones \texttt{UNIQUE} (técnico--zona, caja--número de borne, año--mes).
  \item \textbf{3FN:} los datos de zona, plan, usuario y equipo viven en su propia
  tabla y se referencian por FK. La excepción son los campos \texttt{clientes.equipo\_*},
  que se conservan por compatibilidad con clientes antiguos y ya fueron reemplazados
  por la tabla \texttt{onus}.
\end{itemize}

\begin{table}[htbp]
\centering
\caption{Hallazgos de la validación del modelo y mejoras propuestas.}
\label{tab:ce02-mejoras-bd}
\small
\begin{tabularx}{\textwidth}{L{4cm}XL{3.6cm}}
\toprule
\textbf{Hallazgo} & \textbf{Riesgo} & \textbf{Mejora propuesta} \\
\midrule
RN-01 no está en la base: un cliente puede quedar en dos bornes (comprobado en la prueba). & Inconsistencia de planta si falla la validación de la aplicación. & Índice único parcial \texttt{uq\_borne\_cliente} (probado, ver DDL). \\
Importes en \texttt{float}. & Errores de redondeo en sumas de dinero. & Migrar a \texttt{numeric(10,2)}. \\
Valores por defecto solo en Python (\texttt{default=}); la base no los conoce. & Un \texttt{INSERT} directo falla o queda incompleto (se comprobó con \texttt{cantidad\_tv}). & Agregar \texttt{server\_default} en los modelos. \\
\texttt{password\_pppoe} y \texttt{password\_mikrotik} en texto plano. & Exposición de credenciales si se filtra un respaldo. & Cifrado reversible en reposo (pgcrypto o cifrado en la aplicación) y respaldos cifrados. \\
Campos \texttt{clientes.equipo\_*} duplican \texttt{onus}. & Dos fuentes de verdad sobre el equipo. & Migrar los datos antiguos a \texttt{onus} y eliminarlos. \\
\bottomrule
\end{tabularx}
\end{table}

% ======================================================================
\section{Sección 2: Implementación de la base de datos}
\label{sec:ce02-implementacion-bd}

\subsection{Versionado del esquema con Alembic}

El esquema no se crea a mano: se versiona con \textbf{32 migraciones Alembic}
(14/09/2026 a 18/11/2026) que el contenedor del backend aplica en cada arranque. Cada
cambio del modelo queda como un archivo revisable, lo que permite reconstruir la
base en cualquier servidor.

\begin{table}[htbp]
\centering
\caption{Hitos del historial de migraciones.}
\label{tab:ce02-alembic}
\small
\begin{tabularx}{\textwidth}{L{2cm}L{3cm}X}
\toprule
\textbf{Fecha} & \textbf{Revisión} & \textbf{Cambio} \\
\midrule
14/09/2026 & \texttt{17fad11bb7fa} & Esquema base (\emph{baseline}). \\
17/09/2026 & \texttt{744d882a0b25} & Esquema de empresa de fibra óptica (cajas NAP, bornes). \\
26/09/2026 & \texttt{d879207fee4f} & Mufas y ubicación de zonas. \\
02/10/2026 & \texttt{e99e87341b4b} & Número de operación único en pagos. \\
13/10/2026 & \texttt{32da6de8c64b} & Mensajes masivos de WhatsApp. \\
21/10/2026 & \texttt{e518239e0ea6} & Jornadas y ubicaciones de vendedores. \\
31/10/2026 & \texttt{b8f41d7a0c52} & Clientes empresa (RUC). \\
10/11/2026 & \texttt{a3c7e1f20b48} & Modelo de cuatro roles. \\
15/11/2026 & \texttt{f2a9c81d5e47} & Lista de corte configurable por zona. \\
18/11/2026 & \texttt{c3f1a8d27b45} & Marca de cliente migrado. \\
\bottomrule
\end{tabularx}
\end{table}

\subsection{Entorno de prueba}

Para esta evidencia, el DDL se generó desde los modelos con el dialecto PostgreSQL de
SQLAlchemy~2.0.30 (equivalente al estado \emph{head} de Alembic) y se ejecutó en un
contenedor \texttt{postgres:15-alpine} limpio, junto con datos de prueba
\textbf{anonimizados} y las consultas de verificación. Los tres scripts están en la
carpeta \texttt{sql/} del proyecto y se reproducen con:

\begin{lstlisting}[style=plano]
docker run -d --name wisp_bd -e POSTGRES_PASSWORD=test -e POSTGRES_DB=wisp postgres:15-alpine
docker exec -i wisp_bd psql -U postgres -d wisp < sql/01_ddl.sql
docker exec -i wisp_bd psql -U postgres -d wisp < sql/02_dml.sql
docker exec -i wisp_bd psql -U postgres -d wisp < sql/03_consultas.sql
\end{lstlisting}

\subsection{Script DML (datos de prueba)}

\lstinputlisting[style=sqlwisp,caption={Datos de prueba anonimizados (\texttt{sql/02\_dml.sql}).},label={lst:ce02-dml}]{sql/02_dml.sql}

% ======================================================================
\section{Sección 3: Consultas y programación en base de datos}
\label{sec:ce02-consultas-prog}

\subsection{Consultas SQL y resultados}

Las consultas responden preguntas reales del negocio y prueban las restricciones de
integridad (Q5 y Q6 deben fallar).

\lstinputlisting[style=sqlwisp,caption={Consultas de verificación (\texttt{sql/03\_consultas.sql}).},label={lst:ce02-consultas}]{sql/03_consultas.sql}

\lstinputlisting[style=plano,caption={Resultado de la ejecución en PostgreSQL 15 (18/11/2026).},label={lst:ce02-resultados}]{sql/04_resultado_consultas.txt}

\begin{upeubox}{Lectura de los resultados}
\begin{itemize}[leftmargin=1.2em]
  \item \textbf{Q1} identifica al cliente vencido hace 26 días, candidato a corte.
  \item \textbf{Q2} muestra la ocupación de bornes por caja (5 libres, 2 ocupados, 1 dañado).
  \item \textbf{Q3} calcula ingresos frente a gastos por mes.
  \item \textbf{Q5} y \textbf{Q6} confirman que la base rechaza un número de operación
  repetido (RN-06) y un número de borne duplicado en la misma caja.
  \item \textbf{Q7} lista las zonas que aún no tienen túnel WireGuard.
\end{itemize}
\end{upeubox}


% ----------------------------------------------------------------------
\subsection{Programación en la base de datos: funciones, procedimiento y triggers}
\label{sec:ce02-programacion}

Hoy toda la lógica vive en el backend. Se incorporan a la base las reglas que no deben
depender de que cada pantalla las valide bien, porque su incumplimiento corrompe datos
o deja sin rastro una operación de dinero:

\begin{table}[htbp]
\centering
\caption{Objetos programados en la base de datos.}
\label{tab:ce02-objetos-bd}
\small
\begin{tabularx}{\textwidth}{L{3.6cm}L{2cm}X}
\toprule
\textbf{Objeto} & \textbf{Tipo} & \textbf{Propósito} \\
\midrule
\texttt{sp\_registrar\_pago} & Procedimiento & Valida monto y mes, bloquea la fila del cliente (\texttt{FOR UPDATE}), inserta el pago y, si el cliente estaba suspendido, lo reactiva y lo registra en el historial. Es atómico: o se hace todo, o nada (RN-06, RN-08). \\
\texttt{fn\_dias\_mora} & Función & Días de atraso de un cliente a una fecha; base del corte automático (RN-07). \\
\texttt{trg\_auditoria\_pagos} & Trigger & Copia en \texttt{auditoria\_pagos} cada alta, cambio o baja de un pago, con el usuario de base de datos y el antes/después en JSON (RNF-08). \\
\texttt{trg\_validar\_borne} & Trigger & Impide asignar un cliente a un borne dañado y mantiene el estado \texttt{ocupado}/\texttt{libre} coherente con \texttt{cliente\_id} (RN-01). \\
\texttt{uq\_borne\_cliente} & Índice único parcial & Un cliente ocupa como máximo un borne (RN-01). \\
\bottomrule
\end{tabularx}
\end{table}

\lstinputlisting[style=sqlwisp,caption={Funciones, procedimiento y triggers (\texttt{sql/05\_programacion.sql}).},label={lst:ce02-programacion}]{sql/05_programacion.sql}

\subsubsection*{Pruebas del motor transaccional}

\lstinputlisting[style=sqlwisp,caption={Casos de prueba (\texttt{sql/07\_pruebas\_programacion.sql}).},label={lst:ce02-pruebas-prog}]{sql/07_pruebas_programacion.sql}

\lstinputlisting[style=plano,caption={Resultado de las pruebas en PostgreSQL 15.},label={lst:ce02-res-prog}]{sql/10_resultado_programacion.txt}

\begin{upeubox}{Lectura de las pruebas}
\begin{itemize}[leftmargin=1.2em]
  \item \textbf{P2:} un solo \texttt{CALL} registra el pago, reactiva al cliente suspendido,
  extiende su vencimiento, deja la reactivación en el historial y queda auditado.
  \item \textbf{P3:} un mes mal escrito se rechaza y no queda ningún pago a medias.
  \item \textbf{P4 y P5:} la base impide usar un borne dañado y marca como ocupado el
  borne que recibe un cliente, aunque la aplicación olvide hacerlo.
  \item \textbf{P6 y P7:} el rol de reportes no lee credenciales y la aplicación no puede
  borrar la auditoría (ver \cref{sec:ce02-roles}).
\end{itemize}
\end{upeubox}

% ----------------------------------------------------------------------
\subsection{Planes de ejecución y optimización}
\label{sec:ce02-optimizacion}

Se midieron las dos consultas más frecuentes del sistema ---la ficha del cliente (su
historial de pagos) y la búsqueda de candidatos del corte automático--- sobre una carga
sintética de 50\,000 clientes y 600\,000 pagos, muy por encima del volumen real, para que
la diferencia sea visible. El análisis mostró que \texttt{pagos.cliente\_id}, la columna
por la que más se consulta, \textbf{no tiene índice} en el modelo actual.

\begin{table}[htbp]
\centering
\caption{Efecto de los índices propuestos (EXPLAIN ANALYZE, PostgreSQL 15).}
\label{tab:ce02-indices}
\small
\begin{tabularx}{\textwidth}{XL{2.6cm}L{2.6cm}L{2.4cm}}
\toprule
\textbf{Consulta} & \textbf{Sin índice} & \textbf{Con índice} & \textbf{Mejora} \\
\midrule
Historial de pagos de un cliente (\texttt{ix\_pagos\_cliente\_fecha}) & 21,01 ms (lectura secuencial en paralelo de 600\,000 filas) & 0,086 ms (búsqueda por índice) & unas 240 veces \\
Candidatos a corte (\texttt{ix\_clientes\_estado\_venc}) & 4,24 ms (lectura secuencial de 50\,000 filas) & 0,58 ms (búsqueda por índice) & unas 7 veces \\
\bottomrule
\end{tabularx}
\end{table}

\lstinputlisting[style=plano,caption={Planes de ejecución antes y después de los índices (\texttt{sql/11\_resultado\_rendimiento.txt}).},label={lst:ce02-planes}]{sql/11_resultado_rendimiento.txt}

% ======================================================================
\section{Sección 4: Seguridad y administración}
\label{sec:ce02-seguridad-bd}

\subsection{Usuarios y roles}
\label{sec:ce02-roles}

En el despliegue actual el backend se conecta con un único usuario que es dueño de
todas las tablas. Si ese usuario se filtra, se puede borrar o alterar cualquier dato,
incluida la auditoría. El diseño separa los privilegios en cuatro roles:

\begin{table}[htbp]
\centering
\caption{Roles de la base de datos (mínimo privilegio).}
\label{tab:ce02-roles-bd}
\small
\begin{tabularx}{\textwidth}{L{2.6cm}L{1.4cm}X}
\toprule
\textbf{Rol} & \textbf{Login} & \textbf{Privilegios} \\
\midrule
\texttt{wisp\_owner} & No & Dueño de los objetos; solo lo usan las migraciones de Alembic. \\
\texttt{wisp\_app} & Sí (backend) & \texttt{SELECT/INSERT/UPDATE/DELETE} sobre las tablas, \texttt{EXECUTE} sobre funciones y procedimientos; \textbf{sin} \texttt{UPDATE/DELETE} sobre la auditoría; máximo 30 conexiones. \\
\texttt{wisp\_lectura} & Sí (reportes) & Solo las vistas \texttt{v\_clientes\_reporte} y \texttt{v\_ingresos\_mes}, que no exponen credenciales PPPoE ni PIN. \\
\texttt{wisp\_respaldo} & Sí (copias) & Rol predefinido \texttt{pg\_read\_all\_data}: lee todo, no escribe nada. \\
\texttt{PUBLIC} & --- & Sin privilegios sobre el esquema \texttt{public}. \\
\bottomrule
\end{tabularx}
\end{table}

\lstinputlisting[style=sqlwisp,caption={Roles y privilegios (\texttt{sql/06\_seguridad.sql}). Las contraseñas son de laboratorio.},label={lst:ce02-roles}]{sql/06_seguridad.sql}

\subsection{Auditoría implementada}

La tabla \texttt{historial} ya registra los eventos de negocio del cliente (activación,
corte, pago, cambio de plan u ONU) con su autor. Se agrega \texttt{auditoria\_pagos}, que
registra a nivel de base de datos cada operación sobre los pagos ---el dato económico más
sensible--- con el usuario de conexión y el contenido completo antes y después del
cambio. Como la aplicación no puede modificar ni borrar esa tabla (prueba P7), un
atacante que tome el usuario de la aplicación no puede ocultar sus cambios. Para el resto
de tablas se recomienda la extensión \texttt{pgaudit} cuando la base pase a DB-01 en el
centro de datos (\cref{chap:ce03-dc}).

\subsection{Estrategias de respaldo y recuperación}

\begin{table}[H]
\centering
\caption{Estrategia de respaldo de la base de datos (coherente con el diseño de centro de datos).}
\label{tab:ce02-respaldo}
\small
\begin{tabularx}{\textwidth}{L{3.4cm}L{3.2cm}X}
\toprule
\textbf{Mecanismo} & \textbf{Frecuencia} & \textbf{Detalle} \\
\midrule
\texttt{pg\_dump} en formato personalizado (\texttt{-Fc}) & Diario, 03:00 & Hecho con \texttt{wisp\_respaldo}; hoy ya existe un respaldo diario con retención de 30 días en el servidor. \\
Copia fuera del servidor & Diaria & Al NAS local y cifrada a S3 (\cref{chap:ce03-dc}); hoy es una copia manual semanal. \\
Archivado de WAL & Continuo, cada 15 min & Permite recuperar a un punto en el tiempo (RPO de 15 minutos o menos). \\
Réplica en flujo & Continua & Hacia EC2-DR, para la contingencia (RPO de 5 minutos o menos). \\
Prueba de restauración & Trimestral & Restaurar en una base nueva y comparar conteos y sumas. \\
\bottomrule
\end{tabularx}
\end{table}

Se ejecutó una prueba completa: respaldo con el rol de solo lectura, comprobación de que
ese rol no puede escribir, restauración en una base nueva (simulando la pérdida total) y
verificación de integridad.

\lstinputlisting[style=plano,caption={Prueba de respaldo y restauración (\texttt{sql/13\_resultado\_respaldo.txt}).},label={lst:ce02-respaldo}]{sql/13_resultado_respaldo.txt}

\noindent La base restaurada tiene los mismos 50\,004 clientes, 600\,005 pagos, la misma
suma de montos (29\,400\,245) y los dos triggers. La restauración de unos 4~MB comprimidos
tomó unos 4~segundos, muy por debajo del RTO de 1~hora fijado para la base de datos.

\subsection{Evidencias de monitoreo}

\lstinputlisting[style=sqlwisp,caption={Consultas de monitoreo (\texttt{sql/09\_monitoreo.sql}).},label={lst:ce02-monitoreo}]{sql/09_monitoreo.sql}

\lstinputlisting[style=plano,caption={Resultado del monitoreo en el entorno de prueba.},label={lst:ce02-res-monitoreo}]{sql/12_resultado_monitoreo.txt}

\noindent Estas consultas se programarán en MON-01 (Zabbix) con umbrales: tasa de
aciertos de caché por debajo de 99\,\%, más de 25 conexiones de \texttt{wisp\_app}, tamaño
de la base por encima del 70\,\% del disco de DB-01 y respaldo diario no generado.

\begin{upeubox}{Para pasar a producción}
Los objetos de este capítulo están probados sobre el esquema real, pero todavía no
están en el sistema en producción. Para incorporarlos: (1) una migración Alembic con la
auditoría, los triggers, el procedimiento y los índices; (2) cambiar el
\texttt{DATABASE\_URL} del backend al rol \texttt{wisp\_app}; (3) usar
\texttt{wisp\_respaldo} en el respaldo diario; (4) adaptar \texttt{pagos.py} para llamar a
\texttt{sp\_registrar\_pago}; y (5) migrar los importes a \texttt{numeric(10,2)}
(\cref{tab:ce02-mejoras-bd}).
\end{upeubox}

% ======================================================================
\section{Anexos del entregable}
\label{sec:ce02-e2-anexos}

\begin{table}[H]
\centering
\caption{Scripts y resultados entregados (carpeta \texttt{sql/} del proyecto).}
\label{tab:ce02-scripts}
\small
\begin{tabularx}{\textwidth}{L{5.2cm}X}
\toprule
\texttt{01\_ddl.sql} & DDL completo generado desde los modelos (\cref{lst:ce02-ddl}). \\
\texttt{02\_dml.sql} & Datos de prueba anonimizados. \\
\texttt{03\_consultas.sql}, \texttt{04\_resultado\_consultas.txt} & Consultas de negocio y resultados. \\
\texttt{05\_programacion.sql} & Funciones, procedimiento y triggers. \\
\texttt{06\_seguridad.sql} & Roles, privilegios y vistas. \\
\texttt{07\_pruebas\_programacion.sql}, \texttt{10\_resultado\_programacion.txt} & Pruebas del motor transaccional. \\
\texttt{08\_rendimiento.sql}, \texttt{11\_resultado\_rendimiento.txt} & Planes de ejecución. \\
\texttt{09\_monitoreo.sql}, \texttt{12\_resultado\_monitoreo.txt} & Monitoreo. \\
\texttt{13\_resultado\_respaldo.txt} & Prueba de respaldo y restauración. \\
\bottomrule
\end{tabularx}
\end{table}

\noindent Pendiente: la evidencia de sustentación (exposición breve ante el docente).

\subsection{Script DDL completo}


\lstinputlisting[style=sqlwisp,caption={DDL de la base de datos WISP en PostgreSQL 15 (\texttt{sql/01\_ddl.sql}).},label={lst:ce02-ddl}]{sql/01_ddl.sql}
